\begin{frame}{자주 묻는 질문 (1)}
  \textbf{Q. 정규방정식이 있는데 왜 경사하강법을 배우나요?}
  \vskip0.4em
  \begin{itemize}
    \item 정규방정식은 선형회귀처럼 $J(\boldsymbol{w})$ 가 이차함수라서 미분=0을
      대수적으로 풀 수 있는 \textbf{특수한 경우}에만 존재.
    \item 2.3절 로지스틱회귀부터는 손실(교차 엔트로피)이 $\boldsymbol{w}$ 에 대해
      이차함수가 아니라서 \textbf{애초에 닫힌 형태 해가 없}.
    \item Week 9 신경망은 손실이 \textbf{비볼록}까지.
    \item $\to$ 경사하강법(과 변형들)이 이 학기 전체를 통틀어
      실제로 쓰이는 \textbf{범용 도구}. 정규방정식은 ``정답과 비교해
      경사하강법이 맞게 수렴하는지 검산''하는 용도로 특히 유용.
  \end{itemize}
\end{frame}
